\section{Methods}
Our ideal optimization objective given a training dataset $X, Y = \{x_n, y_n\}_{n=1}^N$ is to maximize recall subject to a minimum precision constraint, as in Eq.~\eqref{eq:max_recall_st_precision}.
While we can easily evaluate this objective for any candidate parameters $\theta$, \emph{training} our model -- that is, finding good values for parameters $\theta$ given a dataset -- is difficult. The key barrier is that the functions involved are based on sums of the flat zero-one function (shown in Fig.~\ref{fig:tp_fp_bounds}) and thus have gradient values equal to zero at almost all $\theta$ values.
We could use gradient-free methods, but these are inefficient when $\theta$ is high-dimensional as it will be in most real settings involving many input variables.
We need to transform this problem so that modern gradient-based learning will be effective.

\citet{ebanScalableLearningNonDecomposable2017} suggest a promising route that leads to an unconstrained objective with non-zero gradients, which we review here and build upon.
First, recognizing that maximizing recall is equivalent to maximizing the true positive count, can rewrite our problem in terms of true and false positive counts:
\begin{align}
\max_{\theta} \text{tpc}(\theta)
\label{eq:max_recall_st_precision_in_terms_of_tpc_fpc},
%\\ \notag
\quad \text{subj. to:}~ 
%\quad  & 
\frac{\text{tpc}(\theta)}
	 {\text{tpc}(\theta) + \text{fpc}(\theta)} \geq \alpha.
\end{align}
Here for simplicity we write the counts as a function of parameters $\theta$ alone. While these counts do depend on the observed training data $X$ and $Y$, we assume this data is known and fixed throughout training.

Next, we can rewrite our problem in an equivalent ``standardized'' form by framing the constraint as a function that must be less than or equal to zero, and minimizing the negative instead of maximizing:
\begin{align}
\min_{\theta} \quad & - \text{tpc}(\theta),
\label{eq:standardized_objective_constrained}
\\ \notag
\text{subj. to:} \quad & 
\underbrace{- \text{tpc}(\theta) + \frac{\alpha}{1-\alpha} \text{fpc}(\theta)}_{g(\theta)} \leq 0.
\end{align}
We emphasize that this transformation is only valid when the sum of $\text{tpc} + \text{fpc}$ is strictly positive, as we need to multiply both sides of the constraint in Eq.~\eqref{eq:max_recall_st_precision_in_terms_of_tpc_fpc} by this sum while preserving the intended inequality.

Using well-established optimization theory \citep{chongChapter23Algorithms2013}, we can transform this to an equivalent unconstrained optimization problem via the penalty method with Lagrange multiplier $\lambda > 0$:
\begin{align}
\min_{\theta} - \text{tpc}(\theta) + \lambda \text{max}(0, g(\theta) )
\label{eq:standardized_objective_UNconstrained}
\end{align}
Here, $g(\theta)$ is the penalty function defined in Eq.~\eqref{eq:standardized_objective_constrained}.
If this penalty function is zero or less than zero, this means the desired constraint is \emph{satisfied}: the precision is at least $\alpha$.

This unconstrained formulation avoids the need to handle a difficult non-linear constraint, but still faces the key problem that the objective as a function of $\theta$ is \emph{flat}. That is, infinitesimal changes in $\theta$ are unlikely to move any example's prediction from one side of the decision boundary to the other, and thus both $\tpc$ and $\fpc$ counts will remain the same for almost all small changes to $\theta$. This flatness is a problem because it means any attempt at gradient-based training will not move the parameters $\theta$ from their original (poor performing) values.
%, obtaining unbiased estimates of the gradient from a minibatch is difficult because the penalty function $g$ is not \emph{decomposable} into a sum over training examples.

\subsection{Previous Bounds using Hinge Loss}
\label{sec:methods_hinge}
\citet{ebanScalableLearningNonDecomposable2017} obtain tractability - informative non-zero gradients at almost all possible $\theta$ parameter values - by defining non-flat \emph{bounds} on the counts using the hinge loss function $h(y, a) = \max(0, 1 - s(y) a)$, where $a \in \mathbb{R}$ is a real value and $s(y)$ is the sign function, which returns $+1$ if $y = 1$ and $-1$ if $y=0$.

First, the false positive count is upper bounded by
\begin{align}
	\text{fpc}^{h}(\theta) = \textstyle \sum_{n : y_n = 0} h(y_n, f_{\theta}(x_n) ).
\end{align}
Similarly, the true positive count is lower bounded by
\begin{align}
	\text{tpc}^{h}(\theta) = \textstyle \sum_{n : y_n = 1} 1 - h(y_n, f_{\theta}(x_n) ).
\end{align}
These bounds can replace the $\text{tpc}$ and $\text{fpc}$ terms in Eq.~\eqref{eq:standardized_objective_UNconstrained}, which is then minimized using a gradient-based solver.
These hinge bounds are visualized in Fig.~\ref{fig:tp_fp_bounds}, where we can observe two key problems.
% with the hinge bounds suggested by \citet{ebanScalableLearningNonDecomposable2017}. 
First, they are \emph{loose} bounds: for just one term in the sum each bound can differ by 1 or more from the ideal value.
Aggregated across all $N_{-}$ or $N_{+}$ terms, the total error can be substantial.
The looseness of the hinge bounds of \citet{ebanScalableLearningNonDecomposable2017} can result in a fundamentally different optimal decision boundary than our proposed tighter sigmoid bounds (introduced later in Sec.~\ref{sec:methods_sigmoid}), as illustrated in Fig~\ref{fig:toy_data_precision_90}.
% and is illustrated visually in our toy 2D example in Figure \ref{fig:toy_data_precision_90}}.
Second, the sum $\fpc^h + \tpc^h$ could be less than zero, as $\tpc^h$ may be negative. This condition would spoil the interpretation of the objective in Eq.~\eqref{eq:standardized_objective_UNconstrained} as equivalent to guaranteeing a specific minimum precision $\alpha$, since we derived it assuming the denominator in Eq.~\eqref{eq:max_recall_st_precision_in_terms_of_tpc_fpc} was strictly positive.


\subsection{New Tighter Sigmoid Bounds}
\label{sec:methods_sigmoid}

We now derive two \emph{families} of functions that bound the zero-one function, one family of upper bounds (for the false positive count) and one family of lower bounds (for the true postive count). These are illustrated in Fig.~\ref{fig:tp_fp_bounds}. 
Both bounds are formed by shifting and scaling the logistic sigmoid function $\sigma(a) = \textstyle \frac{1}{1 + e^{-a}}$.
Concrete bounds can be obtained by fixing tolerance hyperparameters to specific values. By varying tolerances, a user can make the bounds tighter (more accurate but with flatter gradients) or looser (more tractable). Reasonable settings deliver tighter bounds than the hinge bounds of \citet{ebanScalableLearningNonDecomposable2017}.
Furthermore, we'll show that under all conditions our bounds meet the positivity condition required by the denominator in Eq.~\eqref{eq:max_recall_st_precision_in_terms_of_tpc_fpc}.

\textbf{Upper bound on $\fpc$.}
We first seek a non-flat \emph{upper bound} on the zero-one function, denoted $u(a)$. We want this smooth function to meet the following conditions necessary for a tight upper bound:
\begin{align}
\label{eq:fpc_constraint_desiderata}
	u(-\infty)   &\rightarrow 0
\quad && 
	u(-\epsilon) \approx \delta 
\\ \notag 
	u(+\infty) &\rightarrow 1 + \gamma \delta
\quad &&
	u(0) \approx 1 + \delta
\end{align}
where tolerance factors $\epsilon > 0, \delta > 0$ and scaling factor $\gamma \geq 1$ are specified by the user. Large values make the function smoother; small values make the bound tighter. In the limit as $\gamma \rightarrow 1, \epsilon \rightarrow 0, \delta \rightarrow 0$, our function converges to the zero-one function. 

A natural choice for $u(a)$ is to make it a sigmoid whose amplitude is $1+\gamma \delta$ with learnable slope and shift,
\begin{align}
u_{m,b}(a) = (1+\gamma \delta) \sigma( m a + b)
\label{eq:fpc_bound_sigmoid_shift_and_scale},
\end{align}
where the two parameters are slope $m \in \mathbb{R}$ and intercept $b \in \mathbb{R}$. By definition, this class of function satisfies the two limit conditions in Eq.~\eqref{eq:fpc_constraint_desiderata}. To meet the two approximation constraints, we can numerically solve for the optimal parameters that minimize squared error:
\begin{align*}
\hat{m}, \hat{b} =
\argmin_{m \in \mathbb{R}, b \in \mathbb{R}} (\delta - u_{m,b}(-\epsilon) )^2 + (1 + \delta - u_{m,b}(0) )^2.
\end{align*}
Thus, for user-specified tolerance parameters $\gamma,\delta,\epsilon$, we can obtain a tight non-flat differentiable upper bound on the false positive count:
\begin{align}
\fpc^\sigma(\theta) = \textstyle \sum_{n : y_n = 0} (1+\gamma \delta) \sigma( \hat{m} f_{\theta}(x_n) + \hat{b}).
\label{eq:fpc_upper_bound}
\end{align}

If we concretely select a modest setting of our tolerance parameters -- $\gamma=7.00$, $\delta=0.035$, $\epsilon=0.75$ -- we use BFGS to solve the minimization problem and yield slope $\hat{m}=6.85$ and shift $\hat{b}=1.59$. Our choice of tolerance parameters being deemed 'modest' was based on the finding that setting $\gamma, \delta, \epsilon$ yielding tighter bounds did not significantly impact performance on synthetic.


\textbf{Vertically shifting $\tpc$ to enable valid lower bounds without changing the optimal parameter.}
When defining a lower bound for true positive counts, the bound should be strictly positive for all inputs, so that the denominator positivity condition in Eq.~\eqref{eq:max_recall_st_precision_in_terms_of_tpc_fpc} is satisfied.
First, we pick a small positive constant $\gamma \delta$ (where $\gamma > 0, \delta > 0$ are not necessarily the same values used to define $\text{fpc}$ above).
Now, we redefine the true positive count function by adding this small constant to each zero-one function:
%This means each example's contribution to the loss is in $\{\tilde{\gamma} \tilde{\delta}, 1 + \tilde{\gamma} \tilde{\delta\}}$ instead of $\{0, 1\}$. 
\begin{align} \textstyle
    \textnormal{tpc}_{\gamma,\delta}(\theta) = \sum_{n: y_n=1} \gamma \delta + z(f_{\theta}(x_n))
\end{align}
This vertical shift is visualized in Fig.~\ref{fig:tp_fp_bounds}.

Instead of the original unconstrained problem in Eq.~\ref{eq:standardized_objective_UNconstrained}, we then solve a revised optimization problem:
\begin{align}
    \min_{\theta} &- \text{tpc}_{\gamma,\delta}(\theta) + \lambda \text{max}(0, g_{\gamma,\delta}(\theta) ),
    \label{eq:shifted_objective_UNconstrained}
\\ \notag
    g_{\gamma,\delta}(\theta) &= -\textnormal{tpc}_{\gamma,\delta}(\theta) + \textstyle \frac{\alpha}{1-\alpha} \text{fpc}(\theta) + \gamma \delta N_+.
\end{align}
Here the constraint function $g_{\gamma, \delta}$ differs from the original $g$ by substituting $\textnormal{tpc}_{\gamma,\delta}$ and adding constant $\gamma \delta N_{+}$.

\begin{lemma}
Vertically shifting the true positive count does not alter the optimal parameter. Given fixed values of scalars $\lambda, \alpha, \gamma, \delta$, both Eq~\eqref{eq:standardized_objective_UNconstrained} and Eq.~\eqref{eq:shifted_objective_UNconstrained} have the \emph{same} optimal parameter $\theta^*$.
\end{lemma}

\begin{proof}
For all possible $\theta$, we have $g(\theta) = g_{\gamma, \delta}(\theta)$, because inside $g_{\gamma, \theta}$ the constant added in the $-\text{tpc}_{\gamma, \delta}$ term is exactly canceled by the extra additive constant term $\gamma \delta N_+$. Thus, the penalty terms in the two objectives will always yield the same values. The minimization objectives thus differ only by an additive constant ($\text{tpc}$ vs. $\text{tpc}_{\gamma,\delta}$) and will have the same local optima.
\end{proof}
%Instead of using the shifted function in a precision formula, we recommend looking at the unconstrained objective (Eq. 5). Using notation from Sec. 2, let $\textnormal{tpc} = \sum_{n: y_n=1} z(1, f_{\theta}(x_n))$ denote the ideal true positive count and let $\textnormal{tpc}_{\delta}$ be the count where each zero-one function is shifted up by constant $\delta > 0$. By definition, $\textnormal{tpc}_{\delta} = \textnormal{tpc} + \delta N_{+}$. Given fixed values of scalars $\lambda, \alpha, \delta$, the following optimization problems have the \emph{same} optimal parameter $\theta^*$:


\textbf{Lower bound for the shifted $\tpc$.}
We now seek a lower bound $\ell$ of our vertically-shifted zero-one function, meeting the conditions:
\begin{align}
	\ell(-\infty)   &\rightarrow 0
\quad && 
	\ell(0) \approx \delta
\label{eq:tpc_constraint_desiderata}
\\ \notag
	\ell(+\infty) &\rightarrow 1 + \gamma \delta
\quad &&
	\ell(+\tilde{\epsilon}) \approx 1 + \delta
\end{align}
Again tolerance factors $\epsilon > 0, \delta > 0, \gamma \geq 1$ are specified by the user, with similar interpretation as before. 
We emphasize that this lower bound $\ell$ is a distinct function from the upper bound $u$, as shown in Fig.~\ref{fig:tp_fp_bounds}.
Even if we set the tolerance parameters the same for both bounds (which is not necessary but done for simplicity), the optimization desiderata in Eqs.~\eqref{eq:fpc_constraint_desiderata} and \eqref{eq:tpc_constraint_desiderata} are different, yielding distinct functions $\ell$ and $u$.

To obtain a function $\ell$ meeting the goals of Eq.~\eqref{eq:tpc_constraint_desiderata}, we define a sigmoid function with amplitude $1+\gamma \delta$ using the same functional form as Eq.~\eqref{eq:fpc_bound_sigmoid_shift_and_scale}.
Given tolerances $\gamma,\delta,\epsilon$, our lower bound on the vertically-shifted true positive count becomes
\begin{align}
\tpc^\sigma_{\gamma\delta}(\theta) = \textstyle \sum_{n : y_n=1} (1+\gamma\delta) \sigma( \tilde{m} f_{\theta}(x_n) + \tilde{b}).
\label{eq:tpc_lower_bound}
\end{align}
Again, there are two parameters: slope $\tilde{m} \in \mathbb{R}$ and intercept $\tilde{b} \in \mathbb{R}$.
Using the same optimization strategies as for the upper bound above, we solve for the values of $\tilde{m},\tilde{b}$ parameters that best match the desiderata in Eq.~\eqref{eq:tpc_constraint_desiderata}. Concretely fixing tolerances to $\gamma=7.00$, $\delta=0.035$, $\epsilon=0.75$, we obtain $\tilde{m}=6.85$ and $\tilde{b}=-3.54$.

By plugging our lower bound for tpc and our upper bound for fpc into Eq.~\eqref{eq:shifted_objective_UNconstrained}, we achieve a tractable objective amenable to gradient-based training. A strong advantage of using strict \emph{bounds} as we do is the guarantee that if the optimal parameter $\theta^*$ found by optimizing Eq.~\eqref{eq:shifted_objective_UNconstrained} satisfies the constraint penalty (the $g$ term evaluates to less than zero), then the precision will be at least $\alpha$. Simply using differentiable approximations of the zero-one function instead of strict bounds would not provide such a guarantee.

% MCH removed, not necessary
%\textbf{Tightening Bounds Over Epochs} Our sigmoid bound can be tightened over epochs (by varying $\gamma \delta, \epsilon, m, b$), resulting in tighter bounds on the number of true positives and false positives . However, we found empirically on MIMIC, achieves poorer recall at the desired precision constraint. We hypothesize that this is due to the tighter bounds cause gradients to be closer to 0 (due to flatness of the bounds), which hinders training. 


\subsection{Practical Implementation}
Our proposed bounds allow tractable gradient-based optimization of our unconstrained objective in Eq.~\eqref{eq:shifted_objective_UNconstrained}.
Here we discuss practical engineering efforts that make these bounds usable on large real-world datasets.


\textbf{Avoid local optima with many initializations.}
Linear-boundary classifiers like logistic regression (minimizing BCE) or support vector machines (minimizing hinge loss), have \emph{convex} training problems.
However, due to our sigmoid bounds even when our classifier $f_{\theta}$ is a generalized linear model, our objective is \emph{non-convex} and thus sensitive to initialization.
To avoid poor local optima, we take the best of many random initializations (in terms of validation-set performance), where each run is initialized via a ``Glorot'' procedure \citep{glorot2010understanding}.
While keeping the best of many increases runtime if runs are done sequentially, 
separate initializations are easily parallelized across many computers so that overall runtime is no worse than the longest individual run.
Even without full parallelization, we argue our approach is \emph{worthwhile} for the substantial gains in improved recall at the desired precision, especially because training happens much less often than prediction in a deployed alert system. Alternatively, one could develop a homotopy approach \citep{watson1989globally} to navigate the non-convex landscapy and avoid local-optima. However, the success of the homotopy method heavily relies on the choice of the initial simple problem and the path defined using a family of pre-selected parametric function. If the homotopy is not well-designed, it might not provide a useful path, or it might lead to numerical instability.

%In our experiments, we try two kinds of initialization. First, in ``cold start'' runs the weights are initialized via Glorot procedures \citep{glorot2010understanding}. Second, ``warm start'' runs start at small perturbations of the optimal weights found by pretraining on a BCE objective.
%We typically try 25 initializations of each type.

\textbf{Minibatch gradient descent.}
If datasets are small enough, we can optimize our unconstrained objective in Eq.~\eqref{eq:shifted_objective_UNconstrained} using gradients computed from \emph{all} examples.
To scale to larger datasets we use stochastic estimates of the gradient from a minibatch and update parameters via the Adam optimizer~\citep{kingmaAdamMethodStochastic2014}, trying learning rates from 0.0005 to 0.005 and selecting the best performing one on the validation set. 

Using minibatches, however, means that stochastic estimates of the gradient of the $\text{max}(0, g(\theta))$ term in Eq.~\eqref{eq:shifted_objective_UNconstrained} are not unbiased, because the maximum is performed \emph{outside} the sum over all training examples. 
To mitigate this possible bias, we simply select large minibatches (500 or more examples). While not formally unbiased, in practice we find this choice is effective.
We can reach high quality solutions that meet the intended precision when assessed on the entire training set.
To avoid any bias, we could incrementally update per-batch gradients as well as their sum over iterations for a predefined set of batches, an approach known as Stochastic Average Gradient~\citep{schmidtMinimizingFiniteSums2017}.
%In practice, for large datasets (MIMIC and e-ICU) we use the ADAM algorithm with minibatches of size 512, 1024 and the full training set.

% \textbf{Batch size.} In the supplement, we study the sensitivity of performance to batch size. Batches of 512 or 1024 examples are preferred over larger batches; perhaps we might escape local optima better with small batches.

% \textbf{Penalty hyperparameter.} We must choose a value for penalty strength $\lambda > 0$. We found that trying two $\lambda$ values, 1000 and 10000, was sufficient to meet the desired false alarm constraint ($g(\theta) \leq 0$ is satisfied after training). 
% We use these $\lambda$ values in all experiments.

\textbf{Batch size and penalty hyperparameters.} In the supplement, we study the sensitivity of performance to batch size. Batches of 512 or 1024 examples are preferred over larger batches. We found that trying two $\lambda$ values, 1000 and 10000, was sufficient to meet the desired false alarm constraint ($g(\theta) \leq 0$ is satisfied after training). We use these $\lambda$ values in all experiments.

\textbf{Runtime cost analysis.} Our method's training cost will scale \emph{linearly} in the number of examples in each batch. Computing our unconstrained objective in Eq.~\eqref{eq:shifted_objective_UNconstrained} at a specific $\theta$ can be done by evaluating each of the terms $\text{fpc}^{\sigma}$ and $\text{tpc}^{\sigma}_{\gamma,\delta}$ once, which requires just one sigmoidal bound function evaluation ($\ell(\cdot)$ or $u(\cdot)$) for each example (see \eqref{eq:fpc_upper_bound} and \eqref{eq:tpc_lower_bound}).
Computing the \emph{gradient} has the same linear complexity as the loss itself, using well-known complexity results for back-propagation.

\textbf{Chosen objective in practice.}
In pursuit of our ideal but intractable objective in Eq.~\eqref{eq:max_recall_st_precision}, optimizing our tractable unconstrained objective in Eq.~\eqref{eq:shifted_objective_UNconstrained} is guaranteed to produce a valid solution to the original objective if the penalty function $g_{\gamma,\delta}$ is less than or equal to zero. 
In practice, we find that often optimizing Eq.~\eqref{eq:shifted_objective_UNconstrained} is too strict and does not yield a satisfactory result. 
However, another strategy, dropping the additive constant $\gamma \delta N_+$ term in $g_{\gamma,\delta}$ and then post-hoc verifying the desired precision constraint is satisfied, works well. 
We recommend trying both ways and using validation-set performance to select the best.